\documentclass[11pt]{article}
\usepackage[margin=1.1in]{geometry}
\usepackage{amsmath}
\usepackage{graphicx}
\usepackage{booktabs}
\usepackage[colorlinks=true, linkcolor=blue, urlcolor=blue]{hyperref}

\title{How the ray tracing is combined with the Cloudy model:\\
the simulation-to-spectrum pipeline}
\author{pillardisk / transient}
\date{\today}

\newcommand{\nH}{n_{\rm H}}
\newcommand{\rIF}{r_{\rm IF}}
\newcommand{\alphB}{\alpha_{\rm B}}

\begin{document}
\maketitle

\section{Overview}

The pipeline turns one Athena++ dump of the perturbed disk into a
predicted UV--optical spectrum. It has four stages, implemented in
three files under \texttt{transient/}, namely
\texttt{sim\_spectrum.py} (stage 1),
\texttt{sim\_cloudy\_spectrum.py} (stages 2--4), and
\texttt{sim\_full\_spectrum.py} (the thermal and direct
components added in stage~4).

\begin{enumerate}
\item \textbf{Physical fields.} The scale-free simulation cube
      (density, pressure, velocity on a spherical $r,\theta,\phi$ grid)
      is anchored to physical units.
\item \textbf{Ray tracing.} A central ionizing lamp is placed at the
      origin. Along every $(\theta,\phi)$ direction the code marches
      through the simulation's own density cells and finds where the
      lamp's photon budget is exhausted. This defines the
      three-dimensional ionization front, the irradiated ``bowl.''
\item \textbf{Cloudy lookup.} Each ray's irradiated layer is
      characterized by three numbers (ionizing flux, density, column)
      and mapped onto a precomputed Cloudy grid, which returns the
      emitted spectrum of that patch, separately for its illuminated
      face and its shielded back.
\item \textbf{Assembly.} Every patch's spectrum is weighted by its
      area, by which face the observer sees, and by a Lambert
      projection factor, then Doppler shifted by the local gas
      velocity. The sum, plus the thermal light of the photosphere and
      the directly seen lamp, is the predicted spectrum.
\end{enumerate}

The key design choice is that Cloudy is \textbf{not} run per cell.
Cloudy is run once, offline, over a three-parameter family of
irradiated slabs. The ray tracing reduces every line of sight from the
lamp to a member of that family. This keeps the photoionization
physics exact (full radiative transfer within each slab, done by
Cloudy) while the geometry, the shadowing, and the kinematics come
from the simulation.

\section{Stage 1: from simulation to physical fields}

The dump (\texttt{disk.out1.00012.athdf}) is scale free, so three
anchors set the units (all in \texttt{CONFIG}):
\begin{itemize}
\item $r=1$ in code units corresponds to $r_0 = 2$ light days, the
      orbital radius of the disrupted star;
\item the black-hole mass is $M_{\rm BH} = 10^{7}\,M_\odot$, which
      sets the velocity unit $v_0 = \sqrt{GM_{\rm BH}/r_0} \simeq
      5{,}060~{\rm km~s^{-1}}$ and the orbital period at $r=1$
      ($\simeq 740$ days);
\item the azimuthally averaged midplane density at $r=1$ is set to
      $\nH = 10^{10}~{\rm cm^{-3}}$, and the midplane temperature
      anchor there is $3000$~K (the ``rescaled'' temperature mode,
      $T \propto P/\rho$, because the simulation's literal equation of
      state is numerically inflated).
\end{itemize}
Everything downstream (columns, recombination rates, Doppler shifts)
uses these physical fields.

\section{Stage 2: the lamp and the ionization-front march}

A lamp-post ionizing source sits at the origin with photon rate
\begin{equation}
Q_{\rm ion} = 8.5\times10^{52}~{\rm s^{-1}},
\end{equation}
chosen so that the flux arriving at the inner rim is
$\log\varphi \simeq 21$, inside the Cloudy grid. Along each of the
$256\times64$ directions $(\phi,\theta)$ the code marches outward
through the simulation's density cells and integrates the
recombinations per steradian,
\begin{equation}
S(r) \;=\; \int_{r_{\rm in}}^{r} \nH^2\,\alphB\, r'^2\,dr',
\qquad \alphB = 2.59\times10^{-13}~{\rm cm^3\,s^{-1}},
\end{equation}
comparing it with the available photons per steradian,
$Q_{\rm ion}/4\pi$. The ionization front along that ray is the radius
$\rIF(\theta,\phi)$ where $S = Q_{\rm ion}/4\pi$ (photon counting in
Str\"omgren balance, absorption by the gas the lamp actually shines
through). Two outcomes are possible.
\begin{itemize}
\item \textbf{Radiation bounded.} The budget runs out inside the grid.
      The ray has a sharp front, and everything beyond it sits in
      shadow. Dense midplane gas and the arm walls stop the front at
      small radius, so shadowing is automatic; no separate shadow
      computation is needed.
\item \textbf{Matter bounded.} The ray reaches the outer edge of the
      cube with photons to spare (46\% of rays in dump 00012, mostly
      the dilute layers well above the midplane). The whole column is
      ionized and transparent.
\end{itemize}
Note that the photosphere (the $N_{\rm H} = 10^{24}~{\rm cm^{-2}}$
surface) plays no role here. The front emerges from the density field
itself, and the two surfaces are different: the ionization front hugs
the dense structures much more tightly.

\section{Stage 3: reducing each ray to a Cloudy slab}

Cloudy models a plane-parallel slab irradiated on one side. Three
numbers pick the slab for each ray.

\paragraph{Ionizing flux at the front, with foreshortening.}
The front is a tilted surface $r = \rIF(\theta,\phi)$. Its tilt
relative to the incoming radial rays gives the incidence cosine
\begin{equation}
\mu \;=\;
\Bigl[\,1 + \Bigl(\tfrac{\partial \ln \rIF}{\partial\theta}\Bigr)^{2}
        + \Bigl(\tfrac{1}{\sin\theta}
          \tfrac{\partial \ln \rIF}{\partial\phi}\Bigr)^{2}\Bigr]^{-1/2},
\qquad
\varphi \;=\; \frac{Q_{\rm ion}}{4\pi (\rIF r_0)^{2}}\;\mu .
\end{equation}
A wall that faces the lamp (the inner rim, the arm fronts, where
$\rIF$ is nearly constant with angle) has $\mu \simeq 1$ and is hit
hard. A flared surface climbing steeply with angle is grazed,
$\mu \ll 1$. This is the same foreshortening physics as the
irradiation temperature in the pillardisk toy model.

\paragraph{Density.} The slab's hydrogen density is the
recombination-weighted mean over the ray's ionized cells,
$\log \bar{n}_{\rm H} = \sum w_i \log n_{{\rm H},i} / \sum w_i$ with
$w_i = \nH^2 r^2\,dr$ evaluated per cell.

\paragraph{Column density.} A matter-bounded ray hands Cloudy its
true, finite column $N_{\rm H} = \int \nH\,dr$, and Cloudy then
returns a slab that is transparent to ionizing photons. A
radiation-bounded ray is given the grid's maximum column, and Cloudy's
own internal ionization front terminates the slab, so the near
(illuminated) side is fully modeled and the code never has to guess
where Cloudy would place the front.

\section{The precomputed Cloudy grid}

The grid deck (\texttt{transient/cloudy/arm\_column.in}) uses the same
AGN incident spectrum as the older BLR grids (the 15-point
\texttt{interpolate} table), solar abundances, 100~km\,s$^{-1}$
turbulence, and varies three parameters:
\begin{center}
\begin{tabular}{lll}
\toprule
parameter & range (current grid) & range (\texttt{arm\_column2}, queued)\\
\midrule
$\log \varphi$ (ionizing photon flux) & 17 -- 21 & 17 -- 23\\
$\log \nH$ (hydrogen density)         & 9 -- 12  & 6 -- 12\\
$\log N_{\rm H}$ (stopping column)    & 21.5 -- 23.5 & 21.5 -- 24.5\\
\bottomrule
\end{tabular}
\end{center}
For every grid point the \texttt{save continuum} output is split into
four channels, each stored as $\nu F_\nu$ per cm$^2$ of slab surface
on a common wavelength grid (301 -- 11{,}967~\AA):
\begin{itemize}
\item \texttt{refl\_cont}, \texttt{refl\_line}: continuum and lines
      leaving the \textbf{illuminated face} (the ``reflected''
      emission, toward the lamp);
\item \texttt{out\_cont}, \texttt{out\_line}: continuum and lines
      leaving the \textbf{shielded back face} (the ``outward''
      emission, away from the lamp).
\end{itemize}
The four channels are interpolated trilinearly in
$(\log\varphi, \log\nH, \log N_{\rm H})$ in log flux
(\texttt{load\_grid} in \texttt{sim\_cloudy\_spectrum.py}). Rays whose
parameters fall outside the box are clipped to its edge; with the
current grid 62\% of rays clip at the density floor and 68\% at the
column edge, which produces the flat uniform patches in the outer disk
of the line maps. The queued \texttt{arm\_column2} run extends the box
to remove almost all of this clipping.

\section{Face selection and projection}

Resonant lines (Ly$\alpha$, C\,IV) escape preferentially from the
illuminated face, while non-resonant photons leave the shielded back
freely. Cloudy's reflected/outward split encodes exactly this, so the
observer-dependent part of the problem reduces to which face of each
slab patch the observer sees.

The back-face normal of the front surface is
\begin{equation}
\hat{n} \;\propto\; \hat{r}
 - \frac{\partial \ln \rIF}{\partial\theta}\,\hat{\theta}
 - \frac{1}{\sin\theta}\frac{\partial \ln \rIF}{\partial\phi}\,\hat{\phi},
\end{equation}
and with the observer direction $\hat{o}$ (inclination $i$ from the
disk axis, plus an azimuth in the phase sweeps) the blend
\begin{equation}
\cos v = \hat{n}\cdot\hat{o}, \qquad
s_{\rm out} = \tfrac12\bigl[1 + \tanh(\cos v / 0.3)\bigr]
\end{equation}
weights the two faces: $s_{\rm out}\to1$ shows the shielded back
(patch between lamp and observer), $s_{\rm out}\to0$ shows the lit
face, and the smooth blend avoids a hard switch for edge-on patches.
An optically thick patch is also a Lambert emitter, so its flux
carries the projection factor $4|\cos v|$ (the factor 4 makes an
isotropic-equivalent luminosity, so a face-on patch is brighter and an
edge-on patch fainter than the isotropic average, energy conserved
over viewing angles). Matter-bounded (transparent) patches emit both
channels isotropically with no projection factor. At $i=45^\circ$
half the radiation-bounded patches show the observer their lit face
and half their shadow face.

\section{Stage 4: assembling the spectrum}

Each ray's patch has surface area
\begin{equation}
dA \;=\; \frac{(\rIF r_0)^2\, d\Omega}{\mu},
\end{equation}
the solid-angle patch enlarged by $1/\mu$ for the tilt, which together
with $\varphi \propto \mu$ conserves the photon budget per steradian.
Because the observer sees only one face of an opaque disk, only the
observer-side (upper, $\theta \le 90^\circ$) hemisphere enters the
reprocessed sum, matching the one-sided thermal treatment below; the
lower-hemisphere patches have $dA$ set to zero at the source, so the
spectra, both sweeps, and the line maps all inherit the restriction
consistently.
The luminosity spectrum is then
\begin{equation}
L_\lambda \;=\; \sum_{\rm rays} dA \;
\Bigl[ w_{\rm ref}\, F^{\rm refl}_\lambda
     + w_{\rm out}\, F^{\rm out}_\lambda \Bigr]_{\lambda \to \lambda(1 - v_{\rm los}/c)}
\;+\; 4\cos i \; L^{\rm base}_{\lambda,\rm th} \;+\; L_{\lambda,\rm AGN},
\end{equation}
with the three pieces as follows.
\begin{itemize}
\item \textbf{Reprocessed light} (the Cloudy sum). Each patch's entire
      spectrum, continuum and lines alike, is Doppler shifted by the
      recombination-weighted line-of-sight velocity of its ionized
      cells, taken from the simulation velocity field. This is what
      broadens the lines into the observed $\pm3000~{\rm km~s^{-1}}$
      double-horned profiles and smears the Balmer jump; it is not
      applied artificially to lines only.
\item \textbf{Thermal light.} Every disk column radiates as a
      blackbody at the temperature sampled at its own photosphere,
      the depth where the vertical column from the wedge surface
      reaches $N_{\rm H} = 10^{24}~{\rm cm^{-2}}$. This single
      component contains the viscous disk, the TDE-heated region, and
      the arms, because they are all part of the one simulation
      temperature field; the $4\cos i$ factor is the same Lambert
      convention as above.
\item \textbf{Direct AGN light.} The lamp seen directly, using the
      identical \texttt{interpolate} SED that irradiates the Cloudy
      slabs, normalized so that its photon rate above 1~Ryd equals
      $Q_{\rm ion}$. Input SED, reprocessed emission, and direct light
      therefore share one spectral shape and one normalization.
\end{itemize}
The inclination sweep recomputes only face selection, projection, and
Doppler geometry per inclination (the march and the Cloudy lookup are
inclination independent); the phase sweep does the same per observer
azimuth at fixed inclination.

\section{The line maps and the self-consistency check}

The face-on line maps (C\,IV, Mg\,II, H$\alpha$, and the RGB
composite) are \textbf{not} an input to the spectra. Both are
projections of the same per-ray line luminosities. The spectrum sums
them in wavelength; the map paints each ray's line luminosity back
along its ionized cells with the recombination weight
$\nH^2 r^2 dr$ (with the front surface interpolated in $\theta$ so the
rendering shows the continuous bowl rather than the 64 discrete polar
rays).

Numerical check, dump 00012 at $i=45^\circ$ (upper hemisphere).
Summing the H$\alpha$ map gives
$1.83\times10^{40}~{\rm erg~s^{-1}}$, and integrating the line channel
of the $45^\circ$ spectrum over the full Doppler-broadened profile
(6350 -- 6800~\AA) returns the same number, because they are the same
photons, projected spatially in one case and spectrally in the other.
(A narrow window such as 6510 -- 6620~\AA\ underestimates the
spectral side by a factor of a few, since the $\pm3000~{\rm
km~s^{-1}}$ horns carry the flux outside it.)

\section{Approximations worth remembering}

\begin{enumerate}
\item The march uses a single case-B recombination coefficient at
      $10^4$~K; Cloudy then handles the real thermal and ionization
      structure inside each slab, but the location of the front itself
      is photon counting only, with no diffuse-field transport between
      rays.
\item Slabs are characterized by one mean density and one column per
      ray; density gradients along the ionized segment are folded into
      the recombination-weighted mean.
\item First-order Doppler shift only; relativistic beaming
      ($\sim20\%$ on the inner rim) and gravitational redshift
      ($\sim1\%$) are neglected.
\item No patch-on-patch occultation along the observer's line of
      sight, so the sweeps stop at $i = 80^\circ$. For the same reason
      the far-side strip of the lower rim wall that is really visible
      through the central funnel is neglected (the reprocessed sum
      keeps the upper hemisphere only).
\item The lamp is a point at the origin. For
      $M_{\rm BH} = 10^7\,M_\odot$ the gravitational radius is
      $5.7\times10^{-4}$ light days, so any standard lamp height of
      6 -- 20~$r_g$ changes the incidence angles at the simulated radii
      ($\geq 1$ light day) by well under a degree; the physically
      load-bearing part of an elevated lamp, its clear view over the
      inner disk, is present as the assumption that the unsimulated
      hole ($r < 1$ light day) is transparent.
\item 62\% / 68\% of rays currently clip at the Cloudy grid's density
      floor / column edge; superseded once \texttt{arm\_column2} is
      run and extracted (the pipeline adopts the extended cache
      automatically).
\end{enumerate}

\end{document}
